\documentclass[a4paper,11pt]{amsart}
\usepackage[margin=1in]{geometry}
\usepackage[T1]{fontenc}
\usepackage{lmodern,amssymb,mathtools,mathrsfs,esint,xurl}
\usepackage[colorlinks=true,citecolor=blue,linkcolor=blue,urlcolor=blue]{hyperref}
\setlength{\emergencystretch}{2em}
\newtheorem{theorem}{Theorem}[section]
\newtheorem{lemma}[theorem]{Lemma}
\newtheorem{proposition}[theorem]{Proposition}
\newtheorem{corollary}[theorem]{Corollary}
\theoremstyle{remark}\newtheorem{remark}[theorem]{Remark}
\newcommand{\tr}{\operatorname{tr}}
\newcommand{\diag}{\operatorname{diag}}
\newcommand{\PV}{\operatorname{PV}}
\newcommand{\HS}{\mathcal S_2}
\newcommand{\R}{\mathbb R}
\newcommand{\T}{\mathbb T}
\newcommand{\dd}{\,d}
\newcommand{\norm}[1]{\left\lVert#1\right\rVert}
\newcommand{\Cstar}{C_*}
\title{Sharp asymptotics for quadratic Lagrange interpolation}
\date{}
\begin{document}
\begin{abstract}
For distinct interpolation nodes in $[-1,1]$, let $M_n$ denote the minimum of the integral of the sum of the squares of their Lagrange fundamental polynomials.
We prove $M_n=2-C_*/n+o(n^{-1})$, where $C_*=2\Gamma(5/4)^2/\Gamma(3/4)^2$.
The endpoints together with the zeros of $P_{n-2}^{(5/4,5/4)}$ attain this coefficient.
The lower bound follows from exact interpolation identities, compactness of an associated zero-free polynomial, and a weighted Hilbert-transform comparison.
The limiting difference is a nonnegative fractional ground-state integral, with endpoint concentrations retained by lower semicontinuity.
\end{abstract}
\maketitle

\section{Introduction}

For $Y=\{y_1,\ldots,y_n\}\subset[-1,1]$ with distinct elements, write
$\ell_j^Y(x)=\prod_{k\ne j}(x-y_k)/(y_j-y_k)$, and set
\begin{equation}\label{minimum}
 I_n(Y)=\int_{-1}^1\sum_{j=1}^n\ell_j^Y(x)^2\dd x,
 \qquad M_n=\inf_Y I_n(Y).
\end{equation}
The problem of determining $M_n$ goes back to Erd\H{o}s
\cite{erdos}. Fej\'er's solution of the corresponding maximum-norm
problem suggested the Legendre--Gauss--Lobatto points. Szabados
\cite{szabados} disproved their optimality for the integral problem
for every $n>3$. Erd\H{o}s, Szabados, Varma and V\'ertesi
\cite{esvv} proved $2-O((\log n)^2/n)\le M_n\le2-2/(2n-1)$.
Brutman and Toledano \cite{bt} subsequently found numerical evidence
against the proposed first-order coefficient $1$, using optimization
through $n=100$ and extrapolation.

The discussion of problem 1131 \cite{bloom} records a disproof of the
coefficient-one conjecture announced by Price on 26 April 2026, followed
by Sothanaphan's preliminary check on 27 April. An accessible subsequent
proof and Lean development \cite{lean} use the roots of
$T_n-\frac16T_{n-2}$ and establish an eventual coefficient at least
$106/105$. The same discussion records reported exact small-$n$
enclosures, and a preprint \cite{rafik} whose proposed first-order
formula is conditional on an endpoint universality conjecture. Those
results do not supply a sharp unrestricted coefficient. In particular,
the contribution here is the matching lower bound, rather than a new
disproof of coefficient $1$ or a claim that the numerical value had not
been anticipated.

\begin{theorem}\label{main}
For the minimum in \eqref{minimum},
\begin{equation}\label{answer}
 \lim_{n\to\infty}n(2-M_n)
 =\Cstar=2\frac{\Gamma(5/4)^2}{\Gamma(3/4)^2}
 =1.0942198076132386\ldots.
\end{equation}
The arrays $X_n=\{-1,1\}\cup\{x:P_{n-2}^{(5/4,5/4)}(x)=0\}$
satisfy $I_n(X_n)=2-\Cstar/n+o(n^{-1})$.
\end{theorem}

We prove a comparison which is useful independently of minimizers.
Throughout the paper $I=(-1,1)$, $a(x)=1-x^2$, $s(x)=\sqrt{a(x)}$
and $g(x)=a(x)^{1/4}$. Integrals without indicated endpoints are over
$I$. All continuum operators act on $L^2(I,dx)$ unless another measure
is specified. The finite Hilbert transform is
$Hh(x)=\pi^{-1}\PV\int h(t)/(x-t)\dd t$; functions to which $|D|$
is applied are extended by zero to $\R$, and $|D|h=(Hh)'$ on $I$.
The classes $\mathcal S_1$ and $\HS$ denote the trace-class and
Hilbert--Schmidt operators. Multiplication operators are written by
their multiplying functions when this causes no ambiguity.

\begin{theorem}\label{comparison}
If $Y_n$ contains both endpoints and $I_n(Y_n)\le2+C_0/n$ for a fixed
$C_0$, then $\liminf_n n(I_n(Y_n)-I_n(X_n))\ge0$. More precisely,
along a subsequence attaining this lower limit there is a positive
function $f$ such that
\begin{equation}\label{defect}
 \liminf_n n(I_n(Y_n)-I_n(X_n))\ge\mathcal D[f],\qquad
 \mathcal D[f]=\frac1{4\pi}\iint_{I^2}
 \frac{[g(x)f(y)-g(y)f(x)]^2}{f(x)f(y)(x-y)^2}\dd x\dd y<\infty.
\end{equation}
\end{theorem}

The analytic ingredients used below are the classical Hilbert-transform
theorem for Muckenhoupt $A_2$ weights, its reverse-H\"older consequence,
the $L^p$ boundedness of circular conjugation, and standard Hardy-space
boundary theory; see \cite{hmw,gcrf,duren}. On an interval, an $A_2$
weight means a positive function $w$ for which
$\sup_J(\fint_Jw)(\fint_Jw^{-1})<\infty$, with $J$ ranging over its
subintervals. We use the interval versions of these theorems, obtained
by reflection and extension. All problem-specific finite estimates,
compactness statements and endpoint passages are proved below.

\section{Jacobi interpolation and uniform reference estimates}

We use the standard Jacobi normalization $P_m^{(\alpha,\alpha)}(1)
=\Gamma(m+\alpha+1)/[\Gamma(\alpha+1)\Gamma(m+1)]$; see \cite{szego}.
Temporarily fix $1<\alpha<3/2$, put $m=n-2$, $P=P_m^{(\alpha,\alpha)}$,
$q=aP$, $N=m+\alpha+1/2$ and $\beta=3/2-\alpha$. The reference points
are the endpoints and the zeros $x_i$ of $P$; endpoint coordinates are
excluded from sums involving Gaussian weights. Let $\lambda_i$ be the
Gaussian weights for $a^\alpha dx$, and put
\begin{equation}\label{jacobi-constants}
 C_m=2^{2\alpha+1}\frac{\Gamma(m+\alpha+1)^2}
 {\Gamma(m+1)\Gamma(m+2\alpha+1)},\qquad
 A_\alpha=\int a^{\alpha-1},\qquad
 A_m=\sum_i\frac{\lambda_i}{a_i}.
\end{equation}
Here $a_i=a(x_i)$ and $s_i=s(x_i)$. Gaussian exactness and the
Jacobi derivative identity give
$\lambda_i=C_m/[a_iP'(x_i)^2]$ and
$\ell_i^0(x)=q(x)/[a_iP'(x_i)(x-x_i)]$ for interior coordinates.
Define
$W_\alpha(x)=\PV\int a(t)^\alpha/(x-t)\dd t$ and
$F_m(x)=\PV\int a(t)^\alpha P(t)/(x-t)\dd t$.
These normalizations will be retained in every formula involving the
second-kind function $F_m$.

\begin{lemma}\label{jacobi-estimates}
Uniformly on $[-1,1]$,
\begin{equation}\label{PF-bounds}
 |P(x)|\le Cn^{-1/2}(s+n^{-1})^{-\alpha-1/2},\qquad
 |F_m(x)|\le Cn^{-1/2}(s+n^{-1})^{\alpha-1/2}.
\end{equation}
The reference angles have gaps comparable to $n^{-1}$, including the
first and last gaps. The weights $W_i=n\lambda_i a_i^{-\alpha}$ satisfy
$W_i\asymp s_i$ and
\begin{equation}\label{weight-precision}
 |W_i-\pi s_i|\le C/n,\qquad
 W_i=(n/N)\pi s_i+O_K(n^{-2})
\end{equation}
on each compact interior interval $K$ in the second formula.
The full reference square sum is bounded, and
$\int (\ell_\pm^0)^2\le Cn^{-2}$ for the endpoint cardinals.
\end{lemma}

\begin{proof}
The equations for $P$ and $F_m$ are respectively
$aP''-2(\alpha+1)xP'+m(m+2\alpha+1)P=0$ and
$aF_m''+2(\alpha-1)xF_m'+[m(m+2\alpha+1)+2\alpha]F_m=0$.
The second follows by applying the Hilbert transform to the equation
for $a^\alpha P$; that function and its first derivative vanish at the
endpoints. Orthogonality identifies the Wronskian constant:
\begin{equation}\label{wronskian}
 a(PF_m'-P'F_m)+2\alpha xPF_m=C_m.
\end{equation}
Both $u(\theta)=s^{\alpha+1/2}P(\cos\theta)$ and
$v(\theta)=s^{1/2-\alpha}F_m(\cos\theta)$ solve
$y''+[N^2-(\alpha^2-1/4)\csc^2\theta]y=0$.
The center values follow from the Gegenbauer generating function:
if $r_j=(\alpha+1)_j/(2\alpha+1)_j$ and $\lambda=\alpha+1/2$,
$P_{2k}(0)=(-1)^kr_{2k}(\lambda)_k/k!$ and
$P_{2k+1}'(0)=(-1)^k2r_{2k+1}(\lambda)_{k+1}/k!$.
The other center value vanishes by parity. Equation \eqref{wronskian}
then fixes the complementary center data for $F_m$.
Gamma-ratio estimates imply
$|y(\pi/2)|+|y'(\pi/2)|/N\le Cn^{-1/2}$ for both solutions.
Variation of constants and Gronwall give
$|y|+|y'|/N\le Cn^{-1/2}\exp(C|\cot\theta|/N)$.
Thus \eqref{PF-bounds} holds when $s\ge c/n$.

For $0\le\theta\le1/N$, the positive ultraspherical integral
representation gives $|P|\le P(1)$, $|P'|\le P'(1)$ and
$P'(1)/P(1)=m(m+2\alpha+1)/[2(\alpha+1)]$.
Consequently $P\asymp P(1)\asymp n^\alpha$ there. Integrating
$(F_m/(Pa^\alpha))'=C_m/(a^{\alpha+1}P^2)$ from $\cos(1/N)$
proves $|F_m|\le Cn^{-\alpha}$. Reflection treats the other endpoint.
Sturm comparison gives separation at least $c/n$. Beyond a fixed
multiple of $1/n$ the equation has frequency at least $N/\sqrt2$;
comparison there, together with the preceding endpoint layer,
gives the upper gap and first-root bounds.

At a zero, \eqref{wronskian} and \eqref{PF-bounds} bound
$|a_iP'(x_i)|^{-2}\le Cn^{-1}s_i^{2\alpha-1}$, whence
$W_i\le Cs_i$. A local difference-quotient estimate for $u$ gives
$|u'(\theta_i)|\le C\sqrt n$ and therefore the reverse bound,
since $W_i=nC_ms_i/u'(\theta_i)^2$.
For $\mathscr E=u'^2+(N^2-U)u^2$, where
$U=(\alpha^2-1/4)\csc^2\theta$, one has $\mathscr E'=-U'u^2$.
The displayed center formulas and duplication give
$\mathscr E(\pi/2)=NC_m/\pi+O(n^{-1})$; the first-order
gamma-ratio terms cancel. Integration yields
$\mathscr E(\theta)=NC_m/\pi+O((ns^2)^{-1})$.
At a zero this is $u'^2$. Inverting when $ns$ is large, and using
$W_i\le Cs_i$ otherwise, proves \eqref{weight-precision}.

For the remaining assertions, Gaussian partial fractions and
\eqref{wronskian} give the exact identity
\begin{equation}\label{square-identity}
 \sum_j(\ell_j^0)^2-1
 =\frac{B_\alpha P^2-2(\alpha-1)xPF_m}{C_m}
 +\left(\frac{A_m-A_\alpha}{C_m}+\frac1{2P(1)^2}\right)(1+x^2)P^2,
\end{equation}
where
$B_\alpha=A_\alpha(1+x^2-2\alpha)+2(\alpha-1)xW_\alpha$
and $A_\alpha-A_m=C_m/[2\alpha P(1)^2]=O(n^{-2\alpha})$.
One way to check the latter equality is Rodrigues' formula followed
by $m$ integrations by parts in $F_m(1)$.
If $G_b=H(a^b)\pi$, $b=\alpha-1$, then
$B_\alpha=-a^2G_b'$ and $G_b'=O(a^{b-1})$.
Indeed for $0<b<1$, its Pearson equation is
$aG_b'+2bxG_b=2bG_b(1)$, whose odd solution is
$G_b(1)[x+(2b-1)a^b\int_0^x a(t)^{-b}\dd t]$.
Hence $B_\alpha=O(a^\alpha)$, and every term of
\eqref{square-identity} is bounded by \eqref{PF-bounds}.
Finally $\ell_\pm^0=(1\pm x)P/[2P(\pm1)]$ and
$\int P^2=O(n^{2\alpha-2})$, also immediate from
\eqref{PF-bounds}. This proves the endpoint mass estimate.
\end{proof}

\begin{proposition}\label{construction}
For the reference array with parameter $1<\alpha<3/2$,
\begin{equation}\label{cj}
 I_n(X_n^\alpha)=2-\frac{c_J(\alpha)}n+o(n^{-1}),\qquad
 c_J(\alpha)=2\frac{\Gamma(5/2-\alpha)\Gamma(\alpha)}
 {\Gamma(2-\alpha)\Gamma(\alpha-1/2)}.
\end{equation}
Its largest coefficient in this parameter interval is $\Cstar$,
attained at $\alpha=5/4$.
\end{proposition}

\begin{proof}
On compact interior intervals the center-normalized differential
equations in Lemma~\ref{jacobi-estimates} give complementary sine and
cosine expansions, with function errors $O(n^{-3/2})$. In particular
$nP^2/C_m$ converges weakly there to
$1/(2\pi a^{\alpha+1/2})$, while $n\int xPF_m\to0$.
For the last assertion, in angular coordinates the integrand is
$n\cos\theta\,u(\theta)v(\theta)$: it is uniformly bounded outside
the $1/n$ endpoint layers, those layers contribute $O(n^{-1})$,
and interior oscillations have zero mean. The last term in
\eqref{square-identity} has integral $O(n^{-2})$.
Also
$B_\alpha=-a[A_\alpha+(\alpha-1)W_\alpha'/\alpha]$,
with $W_\alpha'$ bounded. Thus the endpoint majorant for
$n|B_\alpha|P^2dx$ is $C\theta^{2-2\alpha}d\theta$ outside the
first layer, and that layer contributes $O(n^{2\alpha-3})$.
It follows that
$c_J=-(2\pi)^{-1}\int B_\alpha a^{-\alpha-1/2}$.

Solving $aW_\alpha'+2\alpha xW_\alpha=2\alpha A_\alpha$ with
$W_\alpha(0)=0$ gives
$W_\alpha'=2\alpha A_\alpha
 {}_2F_1(1,1/2-\alpha;1/2;x^2)$.
Beta integration and Gauss summation \cite[15.4.20]{dlmf} give
\begin{align*}
 c_J&=\frac{B(1/2,\alpha)B(1/2,3/2-\alpha)}{2\pi}
       [1-4(\alpha-1)^2]\\
 &=2\frac{\Gamma(5/2-\alpha)\Gamma(\alpha)}
 {\Gamma(2-\alpha)\Gamma(\alpha-1/2)}.
\end{align*}
For example, the hypergeometric integral needed here is
$\int a^{1/2-\alpha}{}_2F_1(1,1/2-\alpha;1/2;x^2)
 =2(1-\alpha)B(1/2,3/2-\alpha)$; both beta integrals are convergent.
Writing $p=\alpha-1/2$, $q=2-\alpha$, one has $p+q=3/2$ and
$c_J=2\Gamma(p+1/2)\Gamma(q+1/2)/[\Gamma(p)\Gamma(q)]$.
Strict concavity of $\log\Gamma(t+1/2)-\log\Gamma(t)$, whose second
derivative is $\psi_1(t+1/2)-\psi_1(t)<0$, proves the assertion.
\end{proof}

The following estimates are the finite comparison input. For the
interior reference mass matrix $B^0_{ij}=\int\ell_i^0\ell_j^0$, put
$\Omega_n=\diag(\lambda_i a_i^{-\alpha})$, $W_n=n\Omega_n$,
$R_n=n(\Omega_n-B^0)$ and $\mathcal M_n=W_n-R_n$. Extend all these
matrices by zero on endpoint coordinates. Inverses of $W_n$ below
always mean the inverse on the interior coordinates and zero on the
two endpoint coordinates.

\begin{lemma}\label{reference-rows}
There are constants $\rho<1$ and $C$, depending only on $\alpha$, such
that $0\le R_n\le\rho W_n$, $\tr R_n\le C$, and, for every $y\in[-1,1]$,
\begin{equation}\label{row-bounds}
 \left|\sum_i\pi s_i\ell_i^0(y)^2-\pi s(y)\right|\le C/n,
 \qquad L(y)R_nL(y)^*\le C/n,
\end{equation}
where $L(y)=(\ell_i^0(y))_i$. With $k=\alpha-1/2$,
\begin{equation}\label{entry-bound}
 |(R_n)_{ij}|\le\frac Cn
 \left(\frac{\min(s_i,s_j)}{\max(s_i,s_j)}\right)^k.
\end{equation}
\end{lemma}

\begin{proof}
For $c>1$ set
$J_m(c)=\int a^\alpha P^2/(c-x)$ and
$\psi_i(c)=1/[a_iP'(x_i)(c-x_i)]$.
The Gaussian resolvent error is
\begin{equation}\label{resolvent}
 e_m(c)=\int\frac{a^\alpha}{c-x}-\sum_i\frac{\lambda_i}{c-x_i}
 =\frac{J_m(c)}{P(c)^2}>0,
\end{equation}
and the same error with a numerator $p(x)^2$, $\deg p<m$, is
$p(c)^2e_m(c)$. Both statements follow by polynomial division and
Gaussian exactness. Applying
$u^r=(\sin\pi r/\pi)\int_0^\infty t^{r-1}u/(u+t)\dd t$
with $r=2-\alpha$ and $u=a(x)$, and using
$(a+t)^{-1}=[(c-x)^{-1}+(c+x)^{-1}]/(2c)$,
$c=\sqrt{1+t}$, gives the exact positive factorization
\begin{equation}\label{factorization}
 R_n=\frac{n\sin\pi(2-\alpha)}\pi\int_1^\infty
 (c^2-1)^{2-\alpha}J_m(c)
 [\psi(c)\psi(c)^*+\psi(-c)\psi(-c)^*]\dd c.
\end{equation}
Positivity also justifies the integral passages by monotone convergence.
Equation \eqref{PF-bounds} implies
$nJ_m(c)\le C/\sqrt{c^2-1}$.
Together with $|a_iP'(x_i)|^{-2}\le Cn^{-1}s_i^{2k}$, this bounds the
same-end part of an entry, for $u=\min(s_i,s_j)$ and
$v=\max(s_i,s_j)$, by
$Cn^{-1}(uv)^k\int_0^1 h^\beta/[(h+u^2)(h+v^2)]\dd h$.
Splitting at $u^2,v^2$ bounds the integral by $Cv^{-2k}$.
The far part is integrable since $2\beta-2<-1$; for opposite-end
indices one denominator stays bounded near either endpoint.
This proves \eqref{entry-bound}, including $R_{ii}\le C/n$.

For completeness, the strict relative bound is a sampling inequality,
not an assumed spectral limit. Disjoint intervals of length $c/n$
about the reference angles have $\sin\theta\asymp s_i$.
The inequality $|h(t)|^2\le(2/|J|)\int_J|h|^2+2|J|\int_J|h'|^2$
and the Legendre spectral inequality
$\int a|p'|^2\le n(n-1)\int|p|^2$, $\deg p<n$, give
$\sum_i s_i|p(x_i)|^2\le Cn\int|p|^2$.
Since $\Omega_i\le Cs_i/n$, applying this to
$p=\sum z_i\ell_i^0$ gives $B^0\ge c\Omega_n$.
Thus $R_n\le(1-c)W_n$.

We give the endpoint details of both row bounds. Put $f_0=a^{-1/2}$,
$V_\alpha=\PV\int a(t)^{\alpha-1/2}/(x-t)\dd t$, and let
$\varepsilon_m$ be the Gaussian error for the divided difference
$(f_0(y)-f_0(t))/(y-t)$. The Stieltjes representation
$f_0(x)=\pi^{-1}\int_1^\infty(c^2-1)^{-1/2}
 [(c-x)^{-1}+(c+x)^{-1}]\dd c$ gives
$\varepsilon_m'\ge0$ and
$\varepsilon_m'\le Cn^{-2\alpha}s^{-3}$.
Indeed $e_m(c)\le e_m(1)/c$ follows by expanding in even monomials,
whose Gaussian errors are nonnegative, and
$e_m(1)=C_m/[2\alpha P(1)^2]$.
Differentiating the cardinal Cauchy transform and using
\eqref{wronskian} now gives exactly
\begin{equation}\label{weighted-row-exact}
 \sum_i\pi s_i\ell_i^0(y)^2
 =\pi s+\frac\pi{C_m}
 [(1-2\alpha)ysPF_m-a^2P^2V_\alpha'-a^2P^2\varepsilon_m'].
\end{equation}
The Pearson equation used above, with exponent $\alpha-1/2$, gives
$|V_\alpha'|\le Ca^{\alpha-3/2}$.
Consequently the first two corrections are $O(n^{-1})$ by
\eqref{PF-bounds}. The last is bounded by
$Cn^{-2\alpha}sP^2\le C/n$, using separately $s\ge n^{-1}$ and
$s\le n^{-1}$. Continuity includes the endpoints.

For the second row bound in \eqref{row-bounds}, set
$b=\alpha-1$, $G_b=\PV\int a^b/(y-t)\dd t$ and
$H_m(c)=\sum_i\lambda_i/[a_i(c-x_i)]$.
Partial fractions give the exact exterior interpolant
\begin{equation}\label{exterior-mode}
 q_c(y):=L(y)\psi(c)
 =\frac{aP[G_b(y)-H_m(c)]-F_m(y)+y(A_\alpha-A_m)P(y)}{C_m(c-y)}.
\end{equation}
Gaussian positivity in the even-power expansion gives
$0\le G_b(c)-H_m(c)\le Cn^{-2b}$ for $c\ge1$.
To verify the rate, differentiate \eqref{resolvent} at $1$:
if $\Lambda=m(m+2\alpha+1)$ and $D_m=A_\alpha-A_m$, then
$e_m'(1)=-D_m[\alpha\Lambda/(\alpha^2-1)+\alpha/(\alpha-1)]$;
the largest error is
$\int a^{\alpha-2}-\sum_i\lambda_i/a_i^2=(D_m-e_m'(1))/2$.
Also $|G_b(y)-G_b(c)|\le C(s^{2b}+(c-1)^b)$ for
$y\ge0$, $1<c<2$, by the same Pearson formula.
The second-kind equation implies
$|F_m'|\le Cn^{2-\alpha}$ when $s\le n^{-1}$: integrate
$(a^{-b}F_m')'=-(\Lambda+2\alpha)a^{-b-1}F_m$
from $s\asymp n^{-1}$, where the bound follows from the transformed
equation. Since $F_m(1)=D_mP(1)$, it follows that
$|F_m-yD_mP|\le Cs^2n^{2-\alpha}$ in that layer.
Writing $h=c-1$, \eqref{exterior-mode} therefore yields
\begin{align}\label{mode-bounds}
 |q_c(y)|&\le Cn^{-1/2}
 \frac{s^{b+1/2}+s^\beta h^b}{h+s^2}&& (s\ge n^{-1},\ 0<h<1),\notag\\
 |q_c(y)|&\le C\frac{s^2}{h+s^2}
       (n^{1-b}+n^{1+b}h^b)&& (s\le n^{-1},\ 0<h<1).
\end{align}
For $q_{-c}$, $y\ge0$, and also for $q_c$ when $c\ge2$, the bounds
are $Cn^{-1/2}/c$ and $Cn^{b-1}/c$ in the respective ranges.
For the first line of \eqref{mode-bounds}, scaling $h=s^2u$ gives
\begin{equation*}
 s^{2b+1}\int_0^1\frac{h^\beta}{(h+s^2)^2}\dd h\le C,
 \qquad s^{2\beta}\int_0^1\frac{h^{\beta+2b}}{(h+s^2)^2}\dd h\le C.
\end{equation*}
For the second line, split at $h=n^{-2}$; its bounds become
$Cn^{1-b}$ below that point and $Cn^{b-1}h^{b-1}$ above it.
Their squared integrals against $h^\beta dh$ are $O(n^{-1})$,
because $b+\beta=1/2$. The opposite-end and far-exterior integrals
are also $O(n^{-1})$, since
$\int_1^\infty(c^2-1)^\beta c^{-2}\dd c<\infty$.
Insertion into \eqref{factorization} proves
$L(y)R_nL(y)^*\le C/n$. Reflection completes all cases.
\end{proof}

\section{Finite comparison and polynomial bounds}

From now on $\alpha=5/4$ and $X_n=X_n^{5/4}$, although estimates
whose proofs work for $1<\alpha<3/2$ retain that parameter.
For another endpoint-containing array put
$A_{ij}=\ell_j^0(y_i)$, $J_n=A^*A-I$, $U_n=A^*-A^{-1}$ and
$V_n=W_n^{1/2}U_n$. The endpoint rows of $A$ are the corresponding
rows of the identity. A change of the cardinal basis gives the following
exact formula, with the indicated estimate for the omitted mass blocks:
\begin{equation}\label{finite-comparison}
 n(I_n(Y)-I_n(X_n))=-\tr(W_nJ_n)+\tr(R_nJ_n)
 +\tr(\mathcal M_nU_nU_n^*)+\epsilon_{n,e},\qquad
 |\epsilon_{n,e}|\le Cn^{-1/2}\norm{V_n}_{\HS}.
\end{equation}
Indeed $U_nU_n^*=J_n+A^{-1}A^{-*}-I$, and the interior mass is
$(W_n-R_n)/n$. For the endpoint estimate, positivity of the full
old mass matrix gives
$(B^0_{IE})^*(B^0_{II})^{-1}B^0_{IE}\le B^0_{EE}$.
Since $B^0_{II}\le W_n/n$, Lemma~\ref{jacobi-estimates} yields
$\norm{W_n^{-1/2}nB^0_{IE}}_{\HS}\le Cn^{-1/2}$.
If $E=A-I$ and $G=A^{-1}-I$, then $G e_\pm=-U_ne_\pm$;
the endpoint diagonal block of $A^{-1}A^{-*}-I$ is zero.
Cauchy--Schwarz proves \eqref{finite-comparison} without an
unweighted inverse bound.

Let $K_n$ be the reproducing kernel of polynomials of degree less than
$n$ in $L^2(dx)$, and let $\chi_n(x)=K_n(x,x)^{-1}$.
The following uniform estimate will be used also near the endpoints:
\begin{equation}\label{christoffel}
 \chi_n(\cos\theta)=\pi\sin\theta/n+O(n^{-2}),
 \qquad 0\le\theta\le\pi.
\end{equation}
Here is an elementary verification. Christoffel--Darboux gives
$2K_n(x,x)=n^2P_n(x)^2+aP_n'(x)^2$, where $P_n(1)=1$.
Put $u=\sqrt{\sin\theta}P_n(\cos\theta)$ and $N_0=n+1/2$.
Then $u''+[N_0^2+(4\sin^2\theta)^{-1}]u=0$, so for
$E=u^2+(u'/N_0)^2$, $|E'|\le E/(4N_0\sin^2\theta)$.
The center binomial formulas give $E(\pi/2)=2/(\pi N_0)+O(N_0^{-2})$.
Integration gives $E=2/(\pi N_0)+O((N_0^2s)^{-1})$ when
$\theta\in[1/n,\pi-1/n]$. Since
$2sK_n=n^2u^2+(u'-\tfrac12\cot\theta\,u)^2$, this implies
$K_n=n/(\pi s)+O(s^{-2})$. Invert when $ns$ exceeds a sufficiently
large constant. In the remaining endpoint range, Markov's inequality
$|P_j(\cos\theta)-1|\le j^2\theta^2/2$ gives $K_n\ge c n^2$
by summing $j\le c'n$; both terms of \eqref{christoffel} are then
$O(n^{-2})$.

\begin{lemma}\label{coercivity}
Every distinct array satisfies $I_n(Y)\ge2-C/n$. If it contains the
endpoints and $I_n(Y)\le2+C_0/n$, then
\begin{equation}\label{finite-bounds}
 \norm{V_n}_{\HS}\le C,\quad
 \pi\sum_i s(y_i)=2n+O(1),\quad
 \mathscr X_n:=I_n(Y)-\sum_i\chi_n(y_i)=O(n^{-1}),\quad
 \norm{B_Y}\le C/n,
\end{equation}
where $(B_Y)_{ij}=\int\ell_i^Y\ell_j^Y$.
\end{lemma}

\begin{proof}
The residual $r_i=\ell_i^Y-\chi_n(y_i)K_n(\cdot,y_i)$ is orthogonal
to the comparison kernel, so
$\int(\ell_i^Y)^2=\chi_n(y_i)+\norm{r_i}_2^2$.
Consequently \eqref{christoffel} gives
$\pi\sum_i s(y_i)\le nI_n(Y)+C$.
By \eqref{row-bounds}, the uniform reference square-sum bound and
\eqref{weight-precision},
\begin{equation}\label{signed-coarse}
 \tr(W_nJ_n)=\pi\sum_i s(y_i)-2n+O(1)
 \le n(I_n(Y)-2)+C.
\end{equation}
The reference sum is $2n/\pi+O(1)$: the transformed Jacobi equation
gives $N\theta_i=\pi i+O(1)$ uniformly, by integrating its phase
equation from $\theta\asymp n^{-1}$, and summing the Lipschitz
function $\sin\theta$ gives the assertion.

By \eqref{row-bounds}, every row of $E=A-I$ has
$E_iR_nE_i^*\le C/n$. Hence
$\norm{R_n^{1/2}E^*}_{\HS}\le C$ and, writing
$\mathcal P=\norm{V_n}_{\HS}^2$,
$\norm{R_n^{1/2}G}_{\HS}\le C+\sqrt{\rho\mathcal P}$.
As $A^{-1}A^{-*}-I=G+G^*+GG^*$,
$|\tr\{R_n(A^{-1}A^{-*}-I)\}|
 \le\rho\mathcal P+C\sqrt{\mathcal P}+C$.
Using this and \eqref{signed-coarse} in the equivalent form of
\eqref{finite-comparison} gives
\begin{equation}\label{coercive-form}
 2n(I_n(Y)-2)\ge n(I_n(X_n)-2)
 +(1-\rho)\mathcal P-C\sqrt{\mathcal P}-C.
\end{equation}
Proposition~\ref{construction} and completion of the square prove
the lower bound and the first assertion of \eqref{finite-bounds}
for endpoint-containing arrays. The finite identity then also bounds
$\tr(W_nJ_n)$ below, yielding the sine sum and the excess bound.

For the global mass estimate let $\Lambda=\diag(\chi_n(y_i))$ and
$G_{ij}=\sqrt{\chi_n(y_i)\chi_n(y_j)}K_n(y_i,y_j)$.
Then $G_{ii}=1$, $B_Y=\Lambda^{1/2}G^{-1}\Lambda^{1/2}$, and
$P=B_Y+\Lambda^{1/2}G\Lambda^{1/2}-2\Lambda\succeq0$ has
$\tr P=\mathscr X_n$. Thus $B_Y\le2\Lambda+\mathscr X_nI\le CI/n$.

To remove the endpoint hypothesis in the lower bound, write
$S=\sum\ell_i^2$. At either boundary outside the node hull the
extreme cardinal has modulus at least one, and therefore
$S(1),S(-1)\ge1$. Under dilation of the nodes about a point $c$
of their hull the derivative at scale one is
$I_n-(1-c)S(1)-(1+c)S(-1)<0$ if $I_n<2$.
This sign persists throughout admissible dilation. Expand until one
endpoint is reached, then expand about that endpoint until the other
is reached. The integral decreases, proving the unrestricted bound.
The case $n=1$ is immediate.
\end{proof}

For later use we record further consequences with their justifications.
Normalize the nodal polynomial $Q_n$ to leading coefficient $2^{n-1}$.
Then
\begin{equation}\label{Q-bounds}
 c\le\norm{Q_n}_2\le C,\qquad
 n\sum_jQ_n'(y_j)^{-2}\le C.
\end{equation}
Indeed $Q_n=T_n-\mathcal I_YT_n$, and the mass bound in
\eqref{finite-bounds} controls the interpolant of the bounded nodal
data $T_n(y_j)$. The lower bound is the least norm for a prescribed
leading coefficient, attained by the Legendre polynomial; its squared
norm tends to $\pi/4$. To prove the second assertion, put
$Z_n=\sum_j\int(x-y_j)^2\ell_j^Y(x)^2\dd x
 =\norm{Q_n}_2^2\sum_jQ_n'(y_j)^{-2}$.
The residual decomposition above and Christoffel--Darboux give
$Z_n\le C/n+8\mathscr X_n$:
the comparison term for a single node is bounded by
$C\chi_n(y_j)^2[p_{n-1}(y_j)^2+p_n(y_j)^2]\le C/n^2$,
where $p_j=\sqrt{j+1/2}P_j$ and
$p_{n-1}(\cos\theta)^2+p_n(\cos\theta)^2\le C/(s+n^{-1})$.
The latter follows from the same differential-equation bound used in
\eqref{christoffel}; in the first layer use $|P_j|\le1$.

Bounded objective implies that the empirical angular measures tend
to $d\theta/\pi$. To see this without a spacing assumption, let $E$
be the orthonormal-polynomial evaluation matrix. Since
$I_n=\tr((E^*E)^{-1})\le C$, the arithmetic--geometric mean inequality
gives $\det(E^*E)\ge(n/C)^n$.
Factoring the determinant into leading coefficients and the
Vandermonde product gives
\begin{equation*}
 \liminf_{n\to\infty}\frac2{n^2}\sum_{i<j}\log|y_i-y_j|
 \ge-\log2.
\end{equation*}
For any weak limit $\mu$ of the empirical measures, truncating the
logarithmic kernel below gives logarithmic energy at least $-\log2$.
The cosine expansion
$\log|\cos\theta-\cos\phi|
 =-\log2-2\sum_{k\ge1}\cos(k\theta)\cos(k\phi)/k$,
read by Abel regularization, shows that the energy is at most
$-\log2$, with equality only if every nonconstant cosine moment
vanishes. Thus the limit is the arcsine measure. This proves the
assertion and uniform convergence of the ordered angular quantiles
to the uniform quantile map.

\begin{lemma}\label{bulk-spacing}
On each compact subinterval of $(0,\pi)$ the ordered new angles
$\phi_i=\arccos y_i$ have two-sided gaps comparable to $n^{-1}$.
Their normalized kernel Gram matrices satisfy
$c_KI\le G_K\le C_KI$ and $\norm{G_K-I}_{\HS}\le C_K$.
\end{lemma}

\begin{proof}
The excess identity is
$\mathscr X_n=\sum_i\chi_n(y_i)[(G^{-1})_{ii}-1]$.
On a compact interior interval $\chi_n\ge c_K/n$, and Schur
complement inversion gives $(G^{-1})_{KK}\ge G_K^{-1}$.
Since $\tr G_K$ equals its dimension, its eigenvalues $\mu_j$ obey
$\sum_j(\mu_j-1)^2/\mu_j\le C_K$.
This gives the two spectral bounds and the Hilbert--Schmidt bound.
The unit kernel vector
$u_n(\theta)=(p_j(\cos\theta))_{j<n}/\sqrt{K_n(\cos\theta,\cos\theta)}$
has derivative norm at most $C_Kn$. The spectral lower bound on a
two-node principal matrix gives
$\norm{u_n(\phi_i)-u_n(\phi_j)}^2\ge2c_K$, proving separation.

For the upper bound, a gap of length $H/n$ with midpoint $t$ in a
smaller fixed interval admits the nonnegative polynomial
$p_t(\cos\theta)=F_m(\theta-t)+F_m(\theta+t)$, where here
$m=\lfloor n/2\rfloor$ and
$F_m(u)=m^{-1}[\sin(mu/2)/\sin(u/2)]^2$ is the Fej\'er kernel.
It has degree less than $n$, $\norm{p_t}_2\ge c\sqrt n$, and
$0\le p_t(\cos\theta)\le Cn/(1+n|\theta-t|)^2$.
Separation inside a larger compact interval and the bound $O(1/n)$
per sample outside it give
$\sum_i p_t(y_i)\le C[n/(1+H)+1]$.
On the other hand, synthesis and
$\norm{\ell_i^Y}_2\le C/\sqrt n$ give
$\norm{p_t}_2\le Cn^{-1/2}\sum_i p_t(y_i)$.
It follows that $c\le C/(1+H)+C/n$, so $H$ is bounded.
\end{proof}

\section{Boundary convergence of the associated polynomials}

Define analytic polynomials in the disk by
\begin{equation}\label{polynomials}
 F_n(z)=2z^nQ_n((z+z^{-1})/2),\qquad
 b_n(z)=F_n(z)-zF_n'(z)/(2n).
\end{equation}
They have constant coefficient one, and
\begin{equation}\label{boundary-pair}
 b_n(e^{i\theta})=e^{in\theta}
 [Q_n(\cos\theta)-isQ_n'(\cos\theta)/n],\qquad
 f_n(x):=|b_n|^2=Q_n^2+aQ_n'^2/n^2.
\end{equation}
All zeros of $F_n$ are on the unit circle, counted with multiplicity.
The function
$H_n=1-zF_n'/(nF_n)
 =(2n)^{-1}\sum_{\xi:F_n(\xi)=0}(\xi+z)/(\xi-z)$
has positive real part in the disk and $b_n=F_n(1+H_n)/2$.
Thus $b_n$ is zero-free there. Its possible boundary zeros are the
simple zeros at $\pm1$ corresponding to endpoint nodes.

We will use the exact reciprocal identity
\begin{equation}\label{reciprocal}
 \int\frac{s}{f_n}\dd x=\pi n\sum_j Q_n'(y_j)^{-2}\le C.
\end{equation}
Here is a derivation, including the endpoint justification. First move
all nodes into the open interval. Since $\deg b_n<2n$ and it is
zero-free on the closed disk, the identity
$Q_n/f_n=\tfrac12[z^n/b_n(z)+z^{-n}/b_n(z^{-1})]$ on the circle
shows by Fourier orthogonality that $Q_n$ is the degree-$n$
orthogonal polynomial for $dx/(sf_n)$. The same Fourier calculation,
or integration of the Cauchy kernel on the circle, gives
\begin{equation}\label{prufer-cauchy}
 \int\frac{Q_n(t)}{z-t}\frac{dt}{s(t)f_n(t)}
 =\frac{\pi w(z)^n}{\sqrt{z^2-1}\,b_n(w(z))},\qquad
 w(z)=z-\sqrt{z^2-1},
\end{equation}
where the square root is asymptotic to $z$ at infinity.
For clarity, expanding the Cauchy kernel uses
$(z-\cos\theta)^{-1}=(1+2\sum_{k\ge1}w^k\cos k\theta)/\sqrt{z^2-1}$;
only the positive frequencies $z^n/b_n(z)$ remain in the integral.
At a zero $y_j$, the second-kind quotient and
\eqref{prufer-cauchy} give the Gaussian weight
$\lambda_j=\pi/[n f_n(y_j)]=\pi n/[a(y_j)Q_n'(y_j)^2]$.
Gaussian exactness on the polynomial $a$ proves \eqref{reciprocal}.
If an endpoint is a node, move it inward by $\varepsilon$ at fixed
$n$. With $t$ its distance from that endpoint, the denominator obeys
$f_{n,\varepsilon}\ge c_n(t+\varepsilon^2)$: the derivative term
supplies $t$, and the polynomial value supplies $\varepsilon^2$
when $t$ is smaller. Hence $s/f_{n,\varepsilon}\le C_nt^{-1/2}$.
Dominated convergence proves the endpoint version. The bound is now
\eqref{Q-bounds}.

The Legendre derivative inequality and \eqref{Q-bounds} imply
$\int_0^\pi s|b_n|^2d\theta\le C$; \eqref{reciprocal} gives
$\int_0^\pi s^2|b_n|^{-2}d\theta\le C$.
Since $\log^p(1+t)\le C_pt$ and $\log\sin\theta$ belongs to every
finite $L^p$, the positive and negative parts of $u_n=\log|b_n|$
are uniformly bounded in every finite angular $L^p$.
The normalized analytic logarithms $\log b_n(0)=0$ therefore have
bounded $H^p$ norms for $1<p<\infty$, by circular conjugation.
Also $(1-z^2)b_n$ is bounded in $H^2$.
After extraction $F_n$ converges locally uniformly: the generating
function gives $|z^jp_j((z+z^{-1})/2)|\le C_r$ on $|z|\le r<1$;
Cauchy--Schwarz in the Legendre coefficients of $Q_n$, followed by
$\sum_{j\le n}r^{2(n-j)}\le C_r$, bounds $F_n$ on each smaller disk.
The derivative term in \eqref{polynomials} tends locally to zero,
so $b_n$ has the same local limit $F$. Hurwitz's theorem gives a
zero-free $F$ with $F(0)=1$. The logarithms converge locally to
$\log F$ and weakly in $H^p$.
In particular
\begin{equation}\label{outer}
 \log F(z)=\frac1{2\pi}\int_0^{2\pi}
 \frac{e^{it}+z}{e^{it}-z}\log|F(e^{it})|\dd t,
 \qquad \fint_\T\log|F|=0.
\end{equation}
Thus $F$ is outer; this also excludes singular logarithmic terms at
the endpoints before any boundary compactness is used.

\begin{lemma}\label{local-a2}
On every compact interior interval the weights $f_n$ and $f_n^{-1}$
have uniformly bounded $A_2$ characteristics and uniformly bounded
$L^{1+\varepsilon}$ norms for some $\varepsilon>0$.
\end{lemma}

\begin{proof}
Write $d_i=Q_n'(y_i)$. Complex sampling at separated interior points
$t_j+ih/n$ gives, for $\deg p<n$,
$\sum_j|p(t_j+ih/n)|^2\le Cn\norm p_2^2$.
Indeed orthogonality of associated Legendre polynomials gives
$\int a^r|p^{(r)}|^2\le(2n)^{2r}\norm p_2^2$.
The point-evaluation inequality on disjoint intervals of length $c/n$
therefore bounds the derivative samples by
$Cn(C_Kn)^{2r}\norm p_2^2$. Taylor expansion and the $\ell^2$
triangle inequality sum these estimates with factors $(h/n)^r/r!$.
The mass bound in \eqref{finite-bounds} implies that the matrix
$C_{ji}=\ell_i^Y(y_j+i\epsilon)$, $\epsilon=h/n$, is bounded on
every interior square block. Since
$|Q_n(y_j+i\epsilon)/(i\epsilon d_j)|^2
 =\prod_{k\ne j}[1+\epsilon^2/(y_j-y_k)^2]$
is bounded above and below on such rows, the matrix
$D_{ji}=i\epsilon d_j/[d_i(y_j-y_i+i\epsilon)]$ is bounded too.
For a consecutive block $J$ and $Y=\diag(y_j)$,
$[Y,D]+i\epsilon D=i\epsilon(d_j/d_i)_{j,i\in J}$.
Taking the rank-one norm and subtracting the midpoint from $Y$ yields
\begin{equation}\label{discrete-a2}
 \left(\sum_{i\in J}|d_i|^2\right)
 \left(\sum_{i\in J}|d_i|^{-2}\right)\le C_K|J|^2,
\end{equation}
because the block diameter is at most $C_K|J|/n$.
In particular neighboring derivative magnitudes have bounded ratios.

On a consecutive cell $[r,t]$, let $\Delta=t-r$ and
$L(x)=\log|Q_n(x)/[(x-r)(x-t)]|$.
Then $|L''|\le C_Kn^2$, while its two endpoint values differ by a
bounded amount by \eqref{discrete-a2}. Thus $|L'|\le C_Kn$ and
$|Q_n/[(x-r)(x-t)]|\asymp |d_i|/\Delta$ there.
Writing $u=(x-r)/\Delta$, the two terms in \eqref{boundary-pair}
are controlled by $u(1-u)$ and
$1-2u+\Delta L'u(1-u)$ respectively. The first is bounded below
away from the cell ends, and the second is bounded below near those
ends. Consequently $f_n\asymp_K |d_i|^2/n^2$ on the entire cell.
Comparison with cellwise constant weights transfers
\eqref{discrete-a2} to continuous $A_2$; partial cells use the
neighboring ratio bound. The weighted $L^1$ bound for $f_n$, Jensen's
inequality and the global logarithmic bounds control the means of
both weights on fixed intervals. The reverse-H\"older theorem for
$A_2$ weights \cite{gcrf} proves the last assertion.
\end{proof}

\begin{lemma}\label{boundary-compactness}
After extraction, $b_n\to F$ and $b_n^{-1}\to F^{-1}$ strongly in
$L^2$ on compact interior arcs. Their normalized logarithms converge
strongly in every finite angular $L^p$ on the whole circle.
If $\gamma_0(\theta)=\tfrac14(\theta-\pi/2)$ and
$\delta=\operatorname{Im}\log F-\gamma_0$, then the step functions
with values $n(\phi_i-\theta_i)$ converge locally strongly in every
finite $L^p$ to $\delta$.
\end{lemma}

\begin{proof}
Christoffel--Darboux and the center-normalized Legendre equation give,
uniformly on every compact interior angular square,
$G_{ij}=\operatorname{sinc}(n(\phi_i-\phi_j))+O_K(n^{-1})$.
This follows directly by substituting the complementary sine and
cosine expansions of $P_n,P_{n-1}$ in the kernel formula; the
difference between $(2\sin((\phi_i-\phi_j)/2))^{-1}$ and
$(\phi_i-\phi_j)^{-1}$ is bounded, and the diagonal is obtained by
differentiation. In particular, for $h_i=n(\phi_{i+1}-\phi_i)$,
Lemma~\ref{bulk-spacing} gives $c\le h_i\le C$ and
$\sum_i\operatorname{sinc}^2h_i\le C$.
Choose a nearest positive integer multiple $\pi k_i$ except for
the bounded number of gaps near zero. Then
$\sum_i|h_i-\pi k_i|^2\le C$ and
$\sum_i|h_i-\pi k_i|=O(\sqrt n)$.
Equidistribution on each fixed interval gives
$\sum_i h_i-\pi\#\{i\}=o(n)$, so
$\sum_i(k_i-1)=o(n)$ and
\begin{equation}\label{gap-density}
 n^{-1}\sum_i|h_i-\pi|^2\longrightarrow0.
\end{equation}
Thus every fixed number of neighboring gaps approaches the $\pi$
spacing outside a set of indices of density tending to zero.

Fix such interior roots $y_n$ and write $s_n=s(y_n)$.
Define $P_n(t)=nQ_n(y_n+s_nt/n)/(s_nQ_n'(y_n))$.
Its scaled roots tend to $\pi j$, $j\in\mathbb Z$, and
$P_n(0)=0$, $P_n'(0)=1$.
The logarithmic curvature after removing the central and next root
is uniformly bounded on that cell. Neighboring derivative ratios,
from \eqref{discrete-a2}, then bound
$c_n=P_n''(0)/2=-\sum_{j\ne0}r_{j,n}^{-1}$.
Separation gives
$\sum_{|r_{j,n}|>R}|r_{j,n}|^{-2}\le C/R+C/n$.
The exact canonical product
$P_n(t)=t e^{c_nt}\prod_{j\ne0}(1-t/r_{j,n})e^{t/r_{j,n}}$
therefore tends, after extraction, to $e^{ct}\sin t$ locally
uniformly with derivatives. Applying \eqref{discrete-a2} on each
fixed block $-M\le j\le M$ gives
$(\sum e^{2c\pi j})(\sum e^{-2c\pi j})\le C(2M+1)^2$.
Letting $M$ increase forces $c=0$.
A diagonal argument in the window size and the tolerance, using
\eqref{gap-density}, shows that this sine approximation holds on
all but a vanishing fraction of the interior cells.

For $\theta=\phi_i+u/n$ in such a cell, the brackets in
\eqref{boundary-pair}, divided by $s(y_i)Q_n'(y_i)/n$, are
$\sin(-u)-i\cos(-u)+o(1)=-ie^{-iu}+o(1)$ uniformly.
Their modulus is bounded below. Hence
$\sigma_n=b_n/\overline{b_n}$ differs in local $L^2$ by $o(1)$
from the step function with values $-e^{2in\phi_i}$.
The full Gram bound also gives
\begin{equation}\label{circle-energy}
 \sum_{i\ne j}\frac{|e^{2in\phi_i}-e^{2in\phi_j}|^2}{(i-j)^2}
 \le C_K.
\end{equation}
Indeed multiply the sinc kernel entries by
$n(\phi_i-\phi_j)/(i-j)$ and use the two gap bounds; the squared
$O(1/n)$ errors sum to $O(1)$.
Piecewise linear interpolation on the uniform index grid is bounded
in $H^{1/2}$ by \eqref{circle-energy}; pairs of adjacent cells use
the adjacent differences, already included in that sum. Its squared
$L^2$ distance from the step function is at most
$Cn^{-1}\sum_i|e^{2in\phi_{i+1}}-e^{2in\phi_i}|^2=o(1)$.
Compactness gives a strong local $L^2$ limit $\sigma$, with
$|\sigma|=1$ and $\sigma\in H^{1/2}_{\rm loc}$.
The quantile maps between the index grid and the actual cells are
uniformly bi-Lipschitz locally and tend uniformly to the identity.
Composition preserves the strong limit, first for continuous test
functions and then by $L^2$ approximation. Thus $\sigma_n\to\sigma$
locally strongly.

Set $B_n=(1-z^2)b_n$. Weak $H^2$ convergence gives weak local $L^2$
convergence of $b_n$ to $F$; $F$ is nonzero almost everywhere by
\eqref{outer}. Passing to $b_n=\sigma_n\overline{b_n}$ shows
$\sigma=F/\overline F$. We spell out the upgrade to strong amplitude
convergence. For a smooth cutoff $\chi$ in an interior arc let
$\tau=(1-e^{2i\theta})\sigma/\overline{1-e^{2i\theta}}$ and
$h_n=B_n-(1-z^2)F$. Reverse-H\"older in Lemma~\ref{local-a2} and
strong convergence of the bounded quotients imply
$\chi h_n=\tau\chi\overline{h_n}+o_{L^2}(1)$.
Here $h_n$ is weakly null in $H^2$ and has zero constant coefficient.
Let $P_+$ be projection onto the nonnegative circular Fourier modes,
and $P_-=I-P_+$. A smooth cutoff extension $\tau_e$ of $\tau$
belongs to $L^\infty\cap H^{1/2}$ globally. Its commutator with
$P_+$ is Hilbert--Schmidt, because its squared matrix sum is
$\sum_k|k||\widehat{\tau_e}(k)|^2$. Therefore
\begin{align*}
 P_-(\chi h_n)&=[P_-,\chi]h_n\longrightarrow0,\\
 P_+(\chi h_n)&=[P_+,\tau_e]\chi\overline{h_n}
 +\tau_e[P_+,\chi]\overline{h_n}+o_{L^2}(1)\longrightarrow0.
\end{align*}
Both commutators are compact and their inputs are weakly null.
This proves strong local $L^2$ convergence of $b_n$.
Convergence in measure, nonvanishing of $F$, and the uniform
$L^{2+\varepsilon}$ bound for $b_n^{-1}$ prove the reciprocal
convergence by uniform integrability. The same reasoning with
$|\log t|^{p+1}\le C_p(t+t^{-1})$ proves local logarithmic convergence
in every finite $L^p$. The global bounds at a larger exponent make
the endpoint logarithmic tails uniformly small. Thus the real parts
of the logarithms converge globally in every finite $L^p$; circular
conjugation gives the same statement for their imaginary parts.

It remains to fix the real displacement rather than its value
modulo $\pi$. Normalize $\gamma_n=\operatorname{Im}\log b_n$ by
$\log b_n(0)=0$. On $(\phi_i,\phi_{i+1})$, with
$q_n(\theta)=Q_n(\cos\theta)$, the positive-real function $H_n$
above gives exactly
\begin{equation}\label{unwrapped}
 \gamma_n(\theta)=n\theta-\pi i+
 \arctan\frac{q_n'(\theta)}{nq_n(\theta)},\qquad
 \gamma_n(\phi_i)=n\phi_i-\pi(i-1/2).
\end{equation}
The arctangent is its principal value; the normalized logarithm
fixes the displayed branch. On the cells with the sine approximation
the continuous argument differs uniformly by $o(1)$ from its left
nodal value. On the other cells that difference is bounded by the
upper gap constant plus $\pi$, so their total contribution tends
to zero in every finite local $L^p$.
For the reference nodes,
$n\theta_i-\pi(i-1/2)=\gamma_0(\theta_i)+O_K(n^{-1})$,
by the root-anchored Jacobi expansion. Subtracting this identity
from \eqref{unwrapped} and transferring quantile grids proves the
last assertion of the lemma.
\end{proof}

We specify the embeddings used for the remaining matrix limits.
Order both arrays in increasing angle, starting with the endpoint
$1$, and conjugate their matrices by
$\epsilon_i=(-1)^{i+1}$, the signs of the corresponding nodal
derivatives. Map the $i$th coordinate to a constant of height
$\sqrt{n/\pi}$ on $((i-1)\pi/n,i\pi/n)$, and then apply the unitary
map $h(\theta)\mapsto h(\arccos x)/\sqrt{s(x)}$ to physical $L^2$.
Extend $A_n,A_n^{-1}$ by the identity off the step space, and positive
finite mass operators by zero. These conventions preserve every
finite trace and norm identity. In particular $D_n=W_n^{1/2}$ tends
strongly to $D=\sqrt\pi\,g$.

We shall need convergence at both nodal grids. If $p_n$ has
trigonometric degree at most $C_1n$ and $t_j$ are $c/n$ separated
in $K\Subset K_1$, then, for every $M\ge3$,
\begin{equation}\label{sampling}
 \frac1n\sum_{t_j\in K}|p_n(t_j)|^2
 \le C\int_{K_1}|p_n|^2
 +C_{K,K_1,M}n^{1-M}\int_\T|p_n|^2.
\end{equation}
To prove it, reproduce $p_n$ with a smooth Fourier cutoff kernel
$V_n$, for which $|V_n(t)|\le C_Mn(1+n\operatorname{dist}(t,2\pi\mathbb Z))^{-M}$
and $\norm{V_n}_1\le C_M$. Weighted Cauchy--Schwarz in the convolution
and separation bound the sum of the resulting decay kernels by a
constant, and by $Cn^{1-M}$ outside $K_1$.
If analytic polynomials $P_n$ are bounded in $H^2$, have degree
$O(n)$ and converge strongly locally in $L^2$, their step samples
therefore converge strongly locally to the same limit: subtract a
fixed polynomial approximation to the weak $H^2$ limit in
\eqref{sampling}, let $n$ increase and then improve the approximation.
Apply this to $(1-z^2)b_n$ and divide by $1-e^{2i\theta}$ on interior
arcs. On both the reference and new grids the samples of $b_n$
converge strongly locally to $F$. In particular, with
$c=\cos\delta$, $p=\sin\delta$,
\begin{equation}\label{sample-identification}
 \operatorname{Step}_Y\frac{|Q_n'(y_i)|}{n}\longrightarrow\frac{|F|}{s},
 \quad \operatorname{Step}_X(\epsilon_iQ_n(x_i))\longrightarrow|F|p,
 \quad \operatorname{Step}_X\frac{\epsilon_iQ_n'(x_i)}n
 \longrightarrow\frac{|F|c}{s}.
\end{equation}
The first formula uses \eqref{boundary-pair} at a zero; the other
two use the reference root phase in \eqref{unwrapped}.

\begin{proposition}\label{inverse}
Put $f=2^{-1/2}|F(e^{i\arccos x})|^2$,
$T=g^{1/2}Hg^{-1/2}$ and $T_f=f^{1/2}Hf^{-1/2}$.
The operators $D_nA_n^{-1}$ converge weakly to a unique bounded
operator $B$. The product $L=T_fD$ has a bounded extension, and
\begin{equation}\label{channels}
 B=cD+pL,\qquad TB=cL-pD.
\end{equation}
\end{proposition}

\begin{proof}
The finite algebra gives
$\norm{A_nD_nh}^2+\norm{A_n^{-*}D_nh}^2
 =2\norm{D_nh}^2+\norm{U_n^*D_nh}^2$.
Thus $A_nD_n$ and $D_nA_n^{-1}$ are uniformly bounded.
For $h\in C_c^\infty(I)$ the cardinal formula gives exactly
\begin{equation}\label{exact-inverse}
 \sum_j\ell_j^Y(x)\frac{Q_n'(y_j)}n h(y_j)
 =\frac{Q_n'(x)}n h(x)
 -\frac{Q_n(x)}n\sum_j\frac{h(x)-h(y_j)}{x-y_j}.
\end{equation}
The divided difference is continuous on the closed square.
Equidistribution implies uniform convergence of its normalized sum
on compact sets to $-H(h/s)$; the Hilbert transform of the arcsine
density vanishes in $I$.
Give the new grid the raw vector
$d_{j,n}=\sqrt{\pi/n}\,Q_n'(y_j)h(y_j)/n$.
In the physical embedding it tends strongly to
$k_h=|F|h/s^{3/2}$, by \eqref{sample-identification}.
Sampling \eqref{exact-inverse} at the old nodes and multiplying by
$D_n$ gives locally strongly
\begin{equation}\label{B-core}
 Bk_h=\sqrt\pi\,|F|\{ch/s+pH(h/s)\}.
\end{equation}
This is a global $L^2$ function: $H(h/s)$ is bounded near the
endpoints and $\int|F|^2dx\le C$ by the weighted amplitude bound.
The span of $k_h$ is dense in $L^2$ since its multiplying function
is positive almost everywhere and locally square integrable.
Uniform boundedness identifies all weak operator subsequential
limits on this dense family, proving uniqueness and convergence.

The second identity needs an endpoint check. Define on
$\mathbb C\setminus[-1,1]$
$\mathcal R(z)=F(w(z))/(1-w(z)^2)^{1/4}$.
Its upper boundary value is $r e^{-i\delta}$, where
$r=\sqrt{f/g}$, and $\mathcal R(\infty)=1$.
Multiply $\mathcal R$ by the Cauchy transform of $h/s$.
Its real and imaginary boundary parts are
$r[cH(h/s)-ph/s]$ and $-r[ch/s+pH(h/s)]$.
There is no endpoint residue: from $(1-z^2)F\in H^2$, the integral
on a physical circle of radius $\varepsilon$ at either endpoint is
bounded by $C\varepsilon^{1/8}\int_{-\pi}^{\pi}
 |\cos(t/2)|^{-1/2}\dd t=o(1)$.
The compactly supported Cauchy datum is bounded there, and the
product decays at infinity. Cauchy's formula therefore gives
\begin{equation}\label{TB-core}
 TBk_h=\sqrt\pi\,|F|\{cH(h/s)-ph/s\}.
\end{equation}
The fixed transform $T$ is bounded: its power exponents are within
the $A_2$ range. Rotating \eqref{B-core} and \eqref{TB-core} yields
$cB-pTB=D$ and $pB+cTB=T_fD$ on the dense family. The latter
defines the bounded extension $L$, and rotating back proves
\eqref{channels}. Only the product $T_fD$ is asserted bounded.
\end{proof}

\section{The global limiting comparison}

Define $T_0=g^{-1/2}Hg^{1/2}=-T^*$ and
$A=c-pT_0=c+pT^*$. These are bounded operators since the fixed
power weights belong to $A_2$. The reference limit is
\begin{equation}\label{R-limit}
 R_n\longrightarrow R=\tfrac12D(I-T^*T)D
 \quad\hbox{in }\mathcal S_1,
 \qquad K:=TD=\sqrt\pi\,g^{1/2}Hg^{1/2},\quad K^*=-K.
\end{equation}
Here and in what follows finite matrices are embedded as specified
above. We give the global convergence details, because trace
convergence on compact interior squares alone would be insufficient.

In \eqref{factorization} let $v_n(\zeta)$ be the embedded, sign-modulated
vector $-\sqrt{C_m}\psi(\zeta)$. This sign makes the interior nodal
derivatives agree with the positive leading coefficient of $Q_n$.
Root derivative asymptotics give
$v_n(\zeta)\to v(\zeta)=s^{\alpha-1}/(\zeta-x)$ strongly, for each fixed $\zeta>1$.
The endpoint tails follow from
$|v_{n,i}(\zeta)|^2\le Cn^{-1}s_i^{2\alpha-1}/(\zeta-x_i)^2$ and the
reference separation. Also
$nJ_m(\zeta)/C_m\to1/(2\sqrt{\zeta^2-1})$ by the interior oscillatory
average of $na^\alpha P^2/C_m$ and its integrable endpoint bound.
Uniformly in $n$,
$\norm{v_n(\zeta)}^2\le C(\zeta-1)^{\alpha-2}$ for $1<\zeta<2$ and
$\norm{v_n(\zeta)}^2\le C\zeta^{-2}$ for $\zeta\ge2$.
The first bound follows by splitting the angular sum at
$\theta\asymp\sqrt{\zeta-1}$, with the first angular point at $c'/n$.
Thus the trace norm of the factorization integrand is bounded by
$C(\zeta-1)^{-1/2}$ near $1$ and $C\zeta^{2\beta-2}$ at infinity.
Rank-one convergence and dominated convergence in $\mathcal S_1$
prove the first assertion of \eqref{R-limit}.

For identification, put $V_\beta(x)=\PV\int a(t)^\beta/(x-t)\dd t$.
Its Pearson equation, or the odd power series at zero, gives
$V_\beta(x)=\cos(\pi\beta)\int_1^\infty(\zeta^2-1)^\beta
 [(\zeta-x)^{-1}-(\zeta+x)^{-1}]\dd \zeta$.
The limiting kernel from \eqref{factorization} is therefore
\begin{equation}\label{R-kernel}
 R(x,y)=\frac{g(x)^{1/2}g(y)^{1/2}}{2\pi}
 \frac{V_{1/4}(x)-V_{1/4}(y)}{x-y}.
\end{equation}
The finite Hilbert-transform identity
$HM_hH+M_h=[M_{Hh},H]$, with zero extension and initially smooth
compact data, identifies this kernel with
$\tfrac12D(I-T^*T)D$. Extension by boundedness gives
\eqref{R-limit}; in particular $R=\tfrac12(D^2+K^2)$.

We also have the weighted convergences
\begin{equation}\label{weighted-R-limit}
 D_n^{-1}R_n\to D^{-1}R\quad\hbox{in }\HS,
 \qquad D_n^{-1}R_nD_n^{-1}\to\tfrac12(I-T^*T)
 \quad\hbox{strongly}.
\end{equation}
To verify the endpoint assertions, \eqref{entry-bound} and $2k>1$
give $\sum_j|R_{n,ij}|^2\le Cs_i/n$.
For the projection $P_\varepsilon$ onto old coordinates with
$s_i\le\varepsilon$, it follows that
$\norm{P_\varepsilon D_n^{-1}R_n}_{\HS}^2\le C(\varepsilon+n^{-1})$.
For the other orientation use
$R_nW_n^{-1}R_n\le\rho R_n$ and $R_{n,ii}\le C/n$ to obtain
$\norm{D_n^{-1}R_nP_\varepsilon}_{\HS}^2\le C(\varepsilon+n^{-1})$.
Interior kernel convergence proves the first statement of
\eqref{weighted-R-limit}. For the second, test first on compactly
supported functions, control the output tails by the preceding
estimates, and use $0\le D_n^{-1}R_nD_n^{-1}\le\rho I$ and density.
These arguments also give $RD^{-1}\in\HS$.

The weighted forward maps satisfy
\begin{equation}\label{forward}
 A_nD_n\longrightarrow AD\quad\hbox{strongly}.
\end{equation}
Here is a dense-family proof which also controls their endpoint rows.
The vectors $\xi_n(\zeta)=D_n^{-1}v_n(\zeta)$ converge strongly to
$\xi(\zeta)=s^{-1/4}/[\sqrt\pi(\zeta-x)]$.
Indeed their discrete squared coordinates are exactly
$n^{-1}a_i^{\alpha-1}/(\zeta-x_i)^2$; separation controls the tails.
The span of $\xi(\zeta)$, $\zeta>1$, is dense, since expansion in $\zeta^{-1}$
reduces orthogonality to all polynomial moments against the positive
multiplier $s^{-1/4}$.
Substitute the common-phase expansions of $P,F_m$ in the exact
interpolant \eqref{exterior-mode}. Lemma~\ref{boundary-compactness}
then gives $A_nv_n(\zeta)\to Av(\zeta)$ locally strongly; its two terms
are $v(\zeta)c$ and $-pT_0v(\zeta)$, by partial fractions in $G_{\alpha-1}$.
This calculation is valid for unbounded displacement values:
the error from replacing $(N/n)\delta_{i,n}$ by $\delta_{i,n}$
has local mean square tending to zero, since the latter step
functions converge locally in $L^2$.
For each fixed $\zeta>1$, \eqref{exterior-mode} and
\eqref{PF-bounds} give
$\sup_y|q_\zeta(y)|\le C_\zeta n^{-1/2}$.
Each endpoint row thus has squared magnitude at most $C_\zeta/n$;
equidistribution bounds the number of rows in an angular strip
of width $\varepsilon$ by $O(n\varepsilon)+o(n)$.
This upgrades the vector convergence to global strong convergence.
Uniform boundedness of $A_nD_n$, proved in
Proposition~\ref{inverse}, and density prove \eqref{forward}.

Consequently the weak Hilbert--Schmidt limit of $V_n$ is
$V=DA^*-B$, and the positive term in \eqref{finite-comparison} obeys
\begin{equation}\label{positive-liminf}
 \liminf_n\tr(\mathcal M_nU_nU_n^*)
 \ge\tfrac12(\norm V_{\HS}^2+\norm{TV}_{\HS}^2).
\end{equation}
In fact $V_n=(A_nD_n)^*-D_nA_n^{-1}$ is bounded in $\HS$,
and its matrix coefficients are identified by \eqref{forward} and
Proposition~\ref{inverse}. The left side is
$\norm{(I-D_n^{-1}R_nD_n^{-1})^{1/2}V_n}_{\HS}^2$;
strong convergence of the positive square roots and weak lower
semicontinuity prove \eqref{positive-liminf}. Thus no endpoint
concentration has been removed from this term.

The compact trace has an actual limit:
\begin{equation}\label{compact-trace}
 \tr(R_nJ_n)\longrightarrow\tr\{R(A^*A-I)\}.
\end{equation}
On a compact exterior-parameter interval this follows from the
strong vector convergence just proved in \eqref{factorization}.
Its sampled parameter tails are uniformly negligible. To check this,
split new rows into a fixed interior interval and its complement.
For interior rows \eqref{exterior-mode} gives
$|q_{\pm \zeta}(y)|\le C_\varepsilon/(\sqrt n\,\zeta)$, $\zeta\ge1$;
the resulting scalar majorant is
$C_\varepsilon(\zeta^2-1)^\beta/\zeta^2$, integrable at both ends.
For each other row the whole factorization contributes at most
$C/n$ by \eqref{row-bounds}; their total limiting contribution is
$O(\varepsilon)$ by equidistribution. First remove the parameter
tails with $\varepsilon$ fixed and then let $\varepsilon\downarrow0$.
The old trace tails have their independent trace-norm majorant.
This proves \eqref{compact-trace} without a bound for $A_n$ itself.

We next identify the signed trace, keeping its real linear term.
Define
\begin{equation}\label{signed-functionals}
 \mathfrak a(\delta)=\int_0^\pi\cos\theta\,\delta(\theta)\dd\theta,
 \quad Z(\delta)=\frac1\pi\int_0^\pi
 [\pi(1-2\alpha)x\sin\delta\cos\delta
 -s^{3-2\alpha}V_\alpha'(x)\sin^2\delta]\dd\theta,
 \quad \alpha=5/4.
\end{equation}
The coefficients in $Z$ are bounded, including the endpoints.
We claim
\begin{equation}\label{signed-limit}
 \limsup_n\tr(W_nJ_n)\le\mathfrak a(\delta)+Z(\delta).
\end{equation}
Put $W_{0,n}=\diag(\pi s_i)$ and
$Z_n(y)=n[\sum_i\pi s_i\ell_i^0(y)^2-\pi s(y)]$.
The first bound in \eqref{row-bounds} gives $\sup|Z_n|\le C$.
On every compact interior interval \eqref{weighted-row-exact}
and the common-phase expansions give
$Z_n(\cos\theta)=\mathcal Z(x,e^{2id_n(\theta)})+o(1)$, where
\begin{equation}\label{Z-kernel}
 \mathcal Z(x,z)=\tfrac\pi2(1-2\alpha)x\operatorname{Im}z
 -\tfrac12s^{3-2\alpha}V_\alpha'(x)(1-\operatorname{Re}z),
 \qquad d_n(\theta)=N\theta-\tfrac\pi2(\alpha-1/2).
\end{equation}
For example $nPF_m/C_m=\sin d_n\cos d_n/s+O_K(n^{-1})$.
Also $nP^2/C_m=\sin^2d_n/(\pi s^{2\alpha+1})+O_K(n^{-1})$;
the Gaussian error term tends to zero uniformly in the interior.
At new nodes Lemma~\ref{boundary-compactness} gives
$e^{2id_n(\phi_i)}\to e^{2i\delta}$ in step $L^2$.
Since $Z_n$ is uniformly bounded and the endpoint row proportion
tends to zero with the strip width, it follows that
$n^{-1}\sum_iZ_n(y_i)\to Z(\delta)$.

Replacing $W_{0,n}$ by $W_n$ costs $o(1)$ in this trace, not just
$O(1)$. To see the needed precision set
$e_i=W_i-(n/N)\pi s_i$. By \eqref{weight-precision},
$|e_i|\le C/n$ globally and $|e_i|\le C_K/n^2$ in the interior.
The row sum $\sum_i e_i\ell_i^0(y)^2$ is $O(1/n)$ globally.
For interior $y$, split old columns into a larger interior interval
and its complement; the former contributes $O_K(n^{-2})$.
The latter does too, since the cardinal quotient and
$\sum_i|q'(x_i)|^{-2}=O(1)$ give an exterior column square sum
$O_K(n^{-1})$. The same endpoint split as above proves
$\sum_{\rm new}\sum_i e_i\ell_i^0(y_j)^2=o(1)$ and
$\sum_{\rm old}e_i=o(1)$.
Finally $\tr(W_{0,n}J_n)=O(1)$ by \eqref{finite-bounds} and
\eqref{row-bounds}, and $n/N=1+O(1/n)$. We have proved
\begin{equation}\label{signed-split}
 \tr(W_nJ_n)=\pi\left[\sum_i s(y_i)-\sum_i s(x_i)\right]
 +n^{-1}\sum_i Z_n(y_i)+o(1).
\end{equation}

For the remaining linear moment let $Q_n^0$ have the reference
zeros and the same leading coefficient as $Q_n$, and put
$\mathcal R_n=Q_n/Q_n^0$. An exact Stieltjes integral, followed by
integration by parts, gives
\begin{equation}\label{transport}
 \pi\left[\sum_i s(y_i)-\sum_i s(x_i)\right]
 =\int_0^\infty\log|\mathcal R_n(\sqrt{1+t^2})
                            \mathcal R_n(-\sqrt{1+t^2})|\dd t.
\end{equation}
Indeed insert
$a^{1/2}=\pi^{-1}\int_0^\infty u^{-1/2}a/(a+u)\dd u$;
the equal node counts cancel its constant term, and the logarithmic
derivative of the product supplies the remaining resolvents.
At the lower endpoint the boundary term is
$\sqrt{c^2-1}\log|\mathcal R_n(c)\mathcal R_n(-c)|\to0$;
at infinity the odd Laurent coefficients cancel.

Only an upper bound on the positive endpoint tail of the integrand
is needed. Write $\Delta=\sqrt{c^2-1}$ and $\varphi=c+\Delta$.
The Legendre integral representation and \eqref{Q-bounds} give
$|Q_n(c)|+|Q_n(-c)|\le C\varphi^n\min(n,\Delta^{-1})$
for $1<c\le2$. The reference has the lower bound
$|Q_n^0(c)|\ge c_\alpha\varphi^n
 \Delta^2/(\Delta+n^{-1})^{2-\beta}$.
For a direct verification of this last bound, write
$Q_n^0=2m!(x^2-1)C_m^{\alpha+1/2}(x)/(\alpha+1/2)_m$
and use its nonnegative coefficient expansion in
$r=\varphi^{-2}$. Restrict the convolution sum to
$j\le\min(m/2,c/(1-r))$. The coefficient ratio is bounded below,
and the partial binomial sum is at least
$c\min(m,(1-r)^{-1})^{\alpha+1/2}$, which is the asserted bound.
Consequently, for
$h_n(c)=\log|\mathcal R_n(c)\mathcal R_n(-c)|$,
\begin{align*}
 [h_n(c)]_+&\le C+2(1+\beta)\log(1/\Delta)
       &&(n^{-1}\le\Delta\le1),\\
 [h_n(c)]_+&\le C+2(1+\beta)\log n
                  +4\log(1/(n\Delta))&&(0<\Delta<n^{-1}).
\end{align*}
Since $c\,dc/\Delta=d\Delta$, the limiting positive tail on
$0<t<\varepsilon$ is at most
$C\varepsilon(1+|\log\varepsilon|)$.
At infinity $|h_n(c)|\le C/c^2$ uniformly: the normalized analytic
polynomials $F_n$ are bounded and zero-free on a fixed small disk,
and the sum of their logarithms at $w$ and $-w$ has vanishing linear
term. Local uniform convergence, reverse Fatou with the preceding
positive majorants, and \eqref{transport} yield
\begin{equation}\label{transport-upper}
 \limsup_n\pi\left[\sum_i s(y_i)-\sum_i s(x_i)\right]
 \le\int_0^\infty\log|\mathcal R(\sqrt{1+t^2})
                              \mathcal R(-\sqrt{1+t^2})|\dd t.
\end{equation}
The normalized logarithm has the Cauchy representation
$\log\mathcal R(z)=\pi^{-1}\int\delta(x)/(z-x)\dd x$.
This follows from its upper boundary imaginary part $-\delta$;
the endpoint circles vanish because $\log F\in H^p$, $p>1$,
and the value at infinity is zero.
The integral in \eqref{transport-upper} is therefore exactly
$\int x\delta(x)/s(x)\dd x=\mathfrak a(\delta)$:
the inner $t$ integral of
$\pi^{-1}[(\sqrt{1+t^2}-x)^{-1}-(\sqrt{1+t^2}+x)^{-1}]$
is $x/s(x)$. Fubini is absolute since $\delta$ is in angular $L^1$.
Together with \eqref{signed-split} this proves \eqref{signed-limit}.

Combining \eqref{finite-comparison}, \eqref{positive-liminf},
\eqref{compact-trace} and \eqref{signed-limit} proves the fully
global reduction
\begin{equation}\label{global-reduction}
 \begin{gathered}
 \liminf_n n(I_n(Y_n)-I_n(X_n))\ge\mathcal E(F),\qquad
 \\
 \mathcal E(F)=-\mathfrak a(\delta)-Z(\delta)
 +\tr\{R(A^*A-I)\}
 +\tfrac12(\norm V_{\HS}^2+\norm{TV}_{\HS}^2).
 \end{gathered}
\end{equation}
Every term here is finite. Possible negative endpoint concentration
in \eqref{transport} improves the inequality; positive concentration
in \eqref{positive-liminf} has been retained.

\section{The fractional comparison identity}

We first prove finiteness of the nonlinear integral; it is not an
assumption in the approximation argument.

\begin{lemma}\label{finite-defect}
For the limit supplied above, $L-K\in\HS$ and
$\mathcal D[f]=\tfrac14\norm{L+L^*}_{\HS}^2<\infty$.
\end{lemma}

\begin{proof}
The fixed operator $E_0=DT-K=[D,K]D^{-1}$ is Hilbert--Schmidt.
Its kernel is
$\pi^{-1/2}\sqrt{g(x)/g(y)}[g(x)-g(y)]/(x-y)$.
To check the endpoints put $u=1-x$, $v=1-y$ and $q=1/4$.
On $u\le v/2$ its square is bounded by $Cu^qv^{q-2}$; in the
opposite orientation it is bounded by $Cu^{3q-2}v^{-q}$.
Integration reduces both to $C\int_0^1t^{2q-1}dt<\infty$.
Comparable distances use the mean value theorem and give the same
radial integral. Opposite endpoints have a bounded denominator
and integrable power factors.

By \eqref{channels}, $V=DTp-pL$ and $KT=2RD^{-1}-D$.
Therefore the two operators
\begin{equation}\label{two-defects}
 Kp-pL=V-E_0p,\qquad
 Kc-cL=TV-2RD^{-1}p
\end{equation}
are Hilbert--Schmidt. Set $r=\sqrt{f/g}$. On separated compact
rectangles the kernel of $L$ is $r(x)K(x,y)/r(y)$ by its defining
formula in Proposition~\ref{inverse}. The squared kernels of the
two operators in \eqref{two-defects} add to
\begin{equation}\label{kernel-lower}
 |K(x,y)|^2[1+t^2-2t\cos(\delta(x)-\delta(y))]
 \ge |K(x,y)|^2(t-1)^2,\qquad t=r(x)/r(y)>0.
\end{equation}
A countable union of separated rectangles covers the complement
of the diagonal. Hence the difference kernel
$K(x,y)(r(x)/r(y)-1)$ is square integrable globally.
It equals $L-K$ as an operator: on compact intervals the local
$A_2$ realization makes the truncated Hilbert formulas converge
on the dense data family of Proposition~\ref{inverse}; their
difference kernels converge in Hilbert--Schmidt norm. The two
limits agree there, and then globally by exhaustion and boundedness.
This also rules out a remaining diagonal multiplication operator.
Since $K^*=-K$, the symmetric part of $L$ has kernel
$[f(x)g(y)-g(x)f(y)]/[\sqrt\pi\sqrt{f(x)f(y)}(x-y)]$.
Its squared Hilbert--Schmidt norm is exactly $4\mathcal D[f]$.
\end{proof}

For the approximation define
\begin{equation}\label{X-space}
 \norm h_X^2=\int\frac{|h|^2}{s}\dd x+\mathfrak e(h),\qquad
 \mathfrak e(h)=\iint_{I^2}
 \frac{g(x)g(y)|h(x)-h(y)|^2}{(x-y)^2}\dd x\dd y.
\end{equation}
We use $h(\cos\theta)$ as an even circular function and write $Jh$
for its odd circular conjugate, restricted to the upper semicircle.
Then $Jh=-sH(h/s)$ and $H(Jh)=h-\fint_0^\pi h(\cos\theta)d\theta$.
These identities follow on Chebyshev polynomials from
$JT_k=sU_{k-1}$ and $H(sU_{k-1})=T_k$, and then by density.

\begin{lemma}\label{weighted-approximation}
Both $J$ and $H$ are bounded on $X$. If $h$ is in angular $L^2$ and
the kernel $Q_h(x,y)=K(x,y)(e^{h(x)-h(y)}-1)$ is Hilbert--Schmidt,
there are real $h_j\in C_c^\infty(I)$ such that
$h_j\to h$ in $X$ and $Q_{h_j}\to Q_h$ in $\HS$.
\end{lemma}

\begin{proof}
Set $x=\tanh t$ and $a_0(t)=\cosh(t)^{-1/2}$.
Since $d\theta=a_0(t)^2dt$, direct substitution gives
\begin{equation}\label{hyperbolic}
 \mathfrak e(h)=\iint_{\R^2}
 \frac{a_0(t)a_0(u)|h(t)-h(u)|^2}{\sinh^2(t-u)}\dd t\dd u,
 \qquad \norm h_X\asymp\norm{a_0h}_{H^{1/2}(\R)}.
\end{equation}
Here is a verification of the equivalence. For $v=a_0h$, expansion
of the square gives the translation-invariant $1/\sinh^2$ form of
$v$, plus $2\int |v(t)|^2 b(t)dt$, where
$b(t)=\PV\int[a_0(u)/a_0(t)-1]/\sinh^2(t-u)\dd u$.
The function $b$ is bounded: subtract the odd linear Taylor term
near $u=t$, and use $a_0(u)/a_0(t)\le e^{|u-t|/2}$ on the tails.
The reverse expansion has the same bound. The translation-invariant
form plus $L^2$ is equivalent to $H^{1/2}$ because its kernel is
comparable to $|t-u|^{-2}$ near zero and has integrable tails.

The convolution operator $\mathscr H$ with principal-value kernel
$1/(\pi\sinh t)$ is bounded on $H^{1/2}(\R)$, since its Fourier
multiplier is $-i\tanh(\pi\xi/2)$.
The exact transformed formulas are
$J=-\mathscr H$ and $H=a_0^{-2}\mathscr H a_0^2$.
After the isometry $h\mapsto a_0h$, they become
$-a_0\mathscr H a_0^{-1}$ and $a_0^{-1}\mathscr H a_0$.
Their differences from $\mp\mathscr H$ have kernels
$\pi^{-1}[a_0(t)/a_0(u)-1]/\sinh(t-u)$, or the inverse ratio.
These kernels and their first output derivatives have uniform
integrable Schur envelopes. At the diagonal the linear zero
removes the singularity, including after differentiation by
cancellation of the leading terms; on the tails the bound is
$Ce^{-|t-u|/2}$. Thus the differences map $L^2$ into $H^1$.
This proves the asserted bounds on $X$ by \eqref{hyperbolic}.

Symmetry and
$z^2\le C[(e^z-1)^2+(e^{-z}-1)^2]$ give
$\mathfrak e(h)\le C\norm{Q_h}_{\HS}^2$.
Clip $h$ at $\pm N$. Clipping preserves the sign and decreases
the magnitude of every difference, so
$|e^{h_N(x)-h_N(y)}-1|\le|e^{h(x)-h(y)}-1|$.
Dominated convergence gives both $Q_{h_N}\to Q_h$ in $\HS$ and
$h_N\to h$ in $X$.
For fixed $N$, multiply by a smooth hyperbolic cutoff $\chi_R$
equal to one on $[-R,R]$ and zero outside $[-R-1,R+1]$.
The capacity estimate
$\mathfrak e(1-\chi_R)+\norm{1-\chi_R}_{L^2(d\theta)}^2
 \le Ce^{-R}$ follows from \eqref{hyperbolic}, since
$a_0(1-\chi_R)$ and its derivative have squared $L^2$ norms
$O(e^{-R})$. In the difference of $(1-\chi_R)h_N$, the part
containing $h_N(x)-h_N(y)$ is a vanishing tail of its finite
energy; the other part is bounded by $N^2\mathfrak e(\chi_R)$.
Thus these cutoffs converge in $X$. Their uniform boundedness makes
the exponential Lipschitz on the relevant difference range and
gives convergence of $Q$ in $\HS$ as well.
Ordinary interior mollification preserves the bound and converges
in $H^{1/2}$ on compact sets, hence in $X$ and in $Q$.
A diagonal choice proves the lemma.
\end{proof}

We now evaluate the finite expression \eqref{global-reduction}.
First suppose $\delta\in C_c^\infty(I)$ is real, put
$\eta=H\delta$, $r=e^\eta$ and $f=gr^2$.
All ratios $f/g$ are smooth and positive on the closed interval.
The normalized Cauchy factor
$\mathcal R(z)=\exp\{\pi^{-1}\int\delta(t)/(z-t)\dd t\}$ has
boundary value $r e^{-i\delta}$. Its inverse times a Cauchy datum
gives, by Plemelj's formula,
$HM_{r^{-1}}(c-pH)=M_{r^{-1}}(p+cH)$.
Consequently, if $A_g=c-pT$ and $C_g=p+cT$, then
\begin{equation}\label{smooth-inverse}
 C_g=T_fA_g,\qquad A_g^{-1}=c+pT_f,\qquad
 A_gD=DA,\qquad DA^{-1}=A_g^{-1}D.
\end{equation}
For the inverse assertion, $cA_g+pC_g=I$ gives a left inverse.
Both $A_g$ and $A$ are Fredholm of index zero: extended by the
identity outside $I$, they differ by trace-class operators from
$c-pH_{\R}$, whose parametrix is $c+pH_{\R}$ modulo smooth
Hilbert-transform commutators. This applies along the homotopy
$t\delta$, starting from the identity. Thus the left inverse of
$A_g$ is its inverse. The bounded identity $A_gD=DA$ then proves
injectivity of $A$, and index zero proves its invertibility.
No unbounded similarity by $D^{-1}$ has been used.

Here are the trace details, including the trace of a commutator
which cannot be suppressed. Let $U=A^*-A^{-1}$, $J=A^*A-I$,
$\mathsf V=A^{-1}$, $\mathsf F=\mathsf V\mathsf V^*-I$ and
$\mathsf M=D^2-R$. All the differences whose traces occur below
are trace class. This can be checked by expanding the smooth
divided-difference kernels in a Fourier basis: two derivatives in
the compact row or column variable make the sum of rank-one
trace norms finite. The separated endpoint pieces have square
integrable power factors. In particular
$\mathsf V=c+pH_{\R}+\mathcal S_1$.
For $k=a^{1-\alpha}$ the diagonal calculation for $J$, using
$HM_hH+M_h=[M_{Hh},H]$, gives
\begin{equation}\label{trace-J}
 \tr(D^2J)=\int(1-2\alpha)\frac{x}{s}cp\dd x
 -\frac1\pi\int k p^2V_\alpha'\dd x=Z(\delta).
\end{equation}
For precision, the cross-term diagonal is
$-[(cp)'+(k'/k)cp]/\pi$ and the quadratic-term diagonal is
$-(H(kp^2))'/(\pi k)$. Integration by parts against $\pi s$
and the self-adjointness of $|D|$ give \eqref{trace-J}.
Since $UU^*=J+\mathsf F$, \eqref{global-reduction} becomes
$\mathcal E=\tr(\mathsf M\mathsf F)-\mathfrak a(\delta)$.

The required correction is
\begin{equation}\label{trace-correction}
 \tr(\mathsf V^*\mathsf M\mathsf V-\mathsf M)
 -\tr\{\mathsf M(\mathsf V\mathsf V^*-I)\}
 =-\mathfrak a(\delta).
\end{equation}
Indeed the difference is the trace of
$[\mathsf V^*,\mathsf M\mathsf V]$.
The trace-class corrections $R$ and
$\mathsf V-(c+pH_{\R})$ change this trace only by commutators
having a trace-class factor, whose traces vanish.
Write $W=\pi s$ and extend it smoothly to the line near the compact
phase support. The change outside that support contributes a
trace-class separated kernel with zero diagonal and zero trace.
For $[c-H_{\R}p,W(c+pH_{\R})]$ the cross-term diagonal integrates to
\begin{equation*}
 \frac1\pi\int[Wp c'+(Wc)'p]\dd x
 =-\frac1\pi\int W\delta'\dd x=-\mathfrak a(\delta).
\end{equation*}
The quadratic term is
$Q=-H_{\R}(Wp^2)H_{\R}-Wp^2$.
It is trace class and anticommutes with $H_{\R}$; conjugation by
the invertible $H_{\R}$ changes $Q$ to $-Q$, so $\tr Q=0$.
This proves \eqref{trace-correction} by actual traces of
differences, rather than cyclically rearranging divergent traces.
It follows that
$\mathcal E=\tr(\mathsf V^*\mathsf M\mathsf V-\mathsf M)$.

The identities
$A_g^*A_g+C_g^*C_g=I+T^*T$ and \eqref{smooth-inverse} now imply
\begin{equation}\label{smooth-energy}
 \mathcal E=\tfrac12\tr\{D(T_f^*T_f-T^*T)D\}
 =\mathcal C[g]-\mathcal C[f],\qquad
 \mathcal C[h]=\tfrac12\int\frac{s}{h}|D|h\dd x.
\end{equation}
To verify the last equality without an infinite trace subtraction,
the kernel identity above gives
$\mathcal C[h]=\tfrac12\tr\{D(I-T_h^*T_h)D\}$.
For $h=g$ this is the trace of $R$. For $h=f$, the difference
$T_f-T$ has a smooth divided-difference kernel times square
integrable endpoint powers, hence is trace class by the same
separated expansion. The diagonal of the resulting difference is
$s|D|h/h$. Endpoint cutoff limits are justified by the power $1/4$.

The fractional ground-state identity evaluates \eqref{smooth-energy}
explicitly. Direct symmetrization of
$|D|h(x)=\pi^{-1}\PV\int_{\R}[h(x)-h(y)]/(x-y)^2\dd y$ gives
\begin{equation}\label{ground-state}
 \mathcal C[g]-\mathcal C[f]
 =\frac1{4\pi}\iint_{I^2}
 \frac{[g(x)f(y)-g(y)f(x)]^2}{f(x)f(y)(x-y)^2}\dd x\dd y
 =\tfrac14\norm{L+L^*}_{\HS}^2.
\end{equation}
In fact the contribution of $\R\setminus I$ to each $\mathcal C[h]$
is $\pi^{-1}\int s^{-1}=1$, so those contributions cancel.
The remaining numerator is
$s(x)f(y)/f(x)+s(y)f(x)/f(y)-2g(x)g(y)$.
All integrals are finite for a smooth positive ratio $f/g$.
This is the usual quadratic fractional ground-state representation,
also commonly expressed as a Picone identity; see \cite{frank}.
Its elementary expansion here fixes its normalization.
For an independent evaluation of the constant,
$\mathcal C[g]=\tfrac12\langle g,|D|g\rangle$ and
$\widehat g(t)=\sqrt\pi\Gamma(5/4)(2/|t|)^{3/4}J_{3/4}(|t|)$,
with $\widehat h(t)=\int e^{-itx}h(x)dx$.
The beta-integral expansion of the exponential proves this transform.
The Weber--Schafheitlin formula \cite[10.22.57]{dlmf} gives
$\int_0^\infty t^{-1/2}J_{3/4}(t)^2dt
 =\sqrt2/\Gamma(3/4)^2$.
Its convergence conditions hold both at zero and infinity.
Plancherel therefore gives
$\mathcal C[g]=2\Gamma(5/4)^2/\Gamma(3/4)^2=\Cstar$,
in agreement with Proposition~\ref{construction}.

Finally return to the possibly rough limit $F$.
Set $\eta=\log r$, so $\eta$ and $\delta$ are the real and imaginary
parts of $\log[F(z)/(1-z^2)^{1/4}]$ and $\delta=J\eta$.
They are in every finite angular $L^p$.
Lemma~\ref{finite-defect} supplies the nonlinear kernel hypothesis
of Lemma~\ref{weighted-approximation}. Choose its approximants
$\eta_j$ and put $\delta_j=J\eta_j$, $L_j=K+Q_{\eta_j}$,
$p_j=\sin\delta_j$, $c_j=\cos\delta_j$.
Bounded conjugation gives $\delta_j\to\delta$ in $X$.
The sine and cosine compositions also converge in $X$: their
fractional differences have uniformly integrable squares by strong
$X$ convergence, and their Lipschitz constants are at most one.
Thus $[K,p_j]\to[K,p]$ in $\HS$, and likewise for $c_j$.
The identity
$V=E_0p+[K,p]-p(L-K)$ shows that $V_j\to V$ and $TV_j\to TV$
in $\HS$. Multiplication by $c_j,p_j$ converges strongly,
as do its adjoints, so $\tr\{R(A_j^*A_j-I)\}$ converges.
The two signed integrals converge by angular $L^2$ convergence.
Every term of the finite expression \eqref{global-reduction} is
therefore continuous in this approximation, as is
$\tfrac14\norm{L_j+L_j^*}_{\HS}^2$.

For each fixed smooth compact modulus $\eta_j$, its conjugate has
the form $\delta_j=s\psi_j$ with $\psi_j\in C^\infty([-1,1])$;
this follows either from the Cauchy formula or its rapidly decreasing
sine series. Cut this phase off in physical endpoint strips of
width $\varepsilon$. The cutoff tends to $\delta_j$ in $X$ by
the capacity estimate in Lemma~\ref{weighted-approximation}.
Its finite Hilbert conjugate tends to
$\eta_j-\fint\eta_j$ in $X$, and also uniformly with error
$O_j(\sqrt\varepsilon)$. To check the uniform estimate, rescale
$H[(1-\chi_\varepsilon)s\psi_j]$ at each endpoint: it is
$\sqrt\varepsilon$ times the Hilbert transform of a uniformly
bounded $C^{1/2}$ function with square-root vanishing at one end
and smooth vanishing at the other. Subtracting its value in the
principal value leaves the integrable kernel $C|u-v|^{-1/2}$;
the boundary logarithm is bounded by the square-root factor.
Uniform boundedness and the exponential Lipschitz estimate give
convergence of the nonlinear Hilbert--Schmidt kernels as well.
Hence the compact-phase calculation \eqref{ground-state} applies
first to this cutoff, then to $(\eta_j,\delta_j)$, and finally,
by the preceding operator continuity, to $(\eta,\delta)$.
We have proved
\begin{equation}\label{final-identity}
 \mathcal E(F)=\mathcal D[f]=\tfrac14\norm{L+L^*}_{\HS}^2\ge0.
\end{equation}
Together with \eqref{global-reduction} this proves
Theorem~\ref{comparison}.

\begin{proof}[Proof of Theorem~\ref{main}]
The upper comparison $M_n\le I_n(X_n)$ is
Proposition~\ref{construction}. For the reverse asymptotic, a
minimizer exists: in an orthonormal polynomial basis the objective
is $\tr((E^*E)^{-1})$, which diverges when two nodes coalesce;
compactness of the closed ordered node simplex then gives attainment.
For large $n$ the minimum is less than $2$. The endpoint expansion
in Lemma~\ref{coercivity} makes a minimizing array contain both
endpoints. Apply Theorem~\ref{comparison} to these arrays to get
$\liminf_n n(M_n-I_n(X_n))\ge0$. This and \eqref{cj} imply
\eqref{answer}.
\end{proof}

\begin{corollary}\label{rigidity}
If endpoint-containing arrays satisfy
$n(I_n(Y_n)-I_n(X_n))\to0$, then their ordered angular displacement
steps $n(\phi_i-\theta_i)$ tend to zero in every finite $L^p$ on
compact interior intervals.
\end{corollary}

\begin{proof}
Every extracted limit has $\mathcal D[f]=0$, so positivity of the
kernel in \eqref{defect} gives $f/g$ constant almost everywhere.
Outer normalization \eqref{outer} fixes
$F=(1-z^2)^{1/4}$ and therefore $\delta=0$.
Lemma~\ref{boundary-compactness} proves the assertion along every
subsequence, which proves it for the full sequence.
\end{proof}

The supplementary directory \path{verification/} contains independent
finite identity checks, the numerical reproduction programs and their
outputs, and instructions for repeating them. In particular, they
check the cardinal mass identity, the rank-one commutator, the
exterior-resolvent factorization, \eqref{reciprocal},
\eqref{exact-inverse}, and a global smooth-phase comparison.
These calculations provide error-detection and reproducibility;
none is a hypothesis or a substitute for an analytic limiting
argument in the proof. The conclusion concerns the first-order
asymptotic and the interior rigidity above, and does not determine
the exact minimizing arrays for finite $n$.

\begin{thebibliography}{99}
\bibitem{erdos} P. Erd\H{o}s, Problems and results on the theory of
interpolation. II, \emph{Acta Math. Acad. Sci. Hungar.} \textbf{12}
(1961), 235--244.
\bibitem{szabados} J. Szabados, On a problem of P. Erd\H{o}s,
\emph{Acta Math. Acad. Sci. Hungar.} \textbf{17} (1966), 155--157.
\bibitem{esvv} P. Erd\H{o}s, J. Szabados, A. K. Varma and P. V\'ertesi,
On an interpolation theoretical extremal problem,
\emph{Studia Sci. Math. Hungar.} \textbf{29} (1994), 55--60.
\bibitem{bt} L. Brutman and D. Toledano, An extremal problem of
Erd\H{o}s in interpolation theory, \emph{Comput. Math. Appl.}
\textbf{34} (1997), no.~12, 37--47,
\url{https://doi.org/10.1016/S0898-1221(97)00232-0}.
\bibitem{bloom} T. F. Bloom, Erd\H{o}s Problem 1131 and discussion,
\url{https://www.erdosproblems.com/1131},
\url{https://www.erdosproblems.com/forum/thread/1131},
accessed 23 October 2026.
\bibitem{lean} \texttt{seanm27lol}, A Lean disproof of Erd\H{o}s
Problem 1131, mathematical exposition and Lean~4 repository,
\url{https://github.com/seanm27lol/erdos-1131-lean},
accessed 23 October 2026. The formalization is not used here.
\bibitem{rafik} Z. Rafik, Asymptotics of Erd\H{o}s's $L^2$ Lagrange
interpolation problem: arcsine distribution and Airy endpoint
universality, preprint (2026),
\url{https://doi.org/10.20944/preprints202601.0875.v1};
revised manuscript linked in \cite{bloom}.
\bibitem{szego} G. Szeg\H{o}, \emph{Orthogonal Polynomials},
fourth ed., Amer. Math. Soc. Colloq. Publ., vol.~23,
American Mathematical Society, Providence, RI, 1975.
\bibitem{hmw} R. Hunt, B. Muckenhoupt and R. Wheeden,
Weighted norm inequalities for the conjugate function and Hilbert
transform, \emph{Trans. Amer. Math. Soc.} \textbf{176} (1973),
227--251.
\bibitem{gcrf} J. Garc\'ia-Cuerva and J. L. Rubio de Francia,
\emph{Weighted Norm Inequalities and Related Topics},
North-Holland Math. Stud., vol.~116, North-Holland, Amsterdam, 1985.
\bibitem{duren} P. L. Duren, \emph{Theory of $H^p$ Spaces},
Pure Appl. Math., vol.~38, Academic Press, New York, 1970.
\bibitem{frank} R. L. Frank and R. Seiringer, Non-linear ground state
representations and sharp Hardy inequalities,
\emph{J. Funct. Anal.} \textbf{255} (2008), 3407--3430.
\bibitem{dlmf} NIST Digital Library of Mathematical Functions,
\url{https://dlmf.nist.gov/}, especially 10.22.57, 15.4.20 and 18.15.4\_5.
\end{thebibliography}
\end{document}
